% Source: https://artofproblemsolving.com/community/c3h3433271_multiple_choice_test
% Stable identity: aops-topic-3433271-q01
(a) Each symbol appears in a specified position in exactly $(n-1)!$ of the $n!$ permutations, so the probability is $\boxed{\dfrac{1}{n}}$. (b) Summing the $n$ indicator expectations gives $n\cdot\dfrac{1}{n}=\boxed1$. Independence is not required for this step; the correctness events are generally dependent. (c) Exactly one of the $n!$ guesses equals the key, so the probability is $\boxed{\dfrac{1}{n!}}$. For $n=26$, the probability in(a) remains $\dfrac{1}{26}$.
